\scanpage{31}

question de savoir si les conditions constituent bel et bien une \unsure{restriction}, effective, sur un autom. ext. du groupe à lacets $\widehat{\Pi}_{0,3}$, commutant à l'action de $\unsure{\mathfrak{S}}_3$, \unsure{risque} fort d'être hors de portée, faute de savoir construire effectivement de tels autom. extérieurs \inserted{(\unsure{on|ou} \illegible{1} des \unsure{élémts} \unsure{non triviaux} de $\unsure{\widehat{\widehat{\mathcal{T}}}|\widehat{\mathcal{T}}}_{0,3}^{!\unsure{*}}$)} d'une autre façon que justement par voie arithmétique, à l'aide de l'opération de $\Gamma_{\mathbb{Q}}$ sur $\widehat{\Pi}_{0,3}$ ! (Il faudrait quand \unsure{même} refaire à l'occasion un effort en ce sens \ldots)

Je me pose aussi la question dans quelle mesure on peut reconstituer, en termes de la théorie des groupes profinis à lacets et des \deleted{\unsure{leurs}} autom. ext. de $\widehat{\Pi}_{0,3}$, ou de \deleted{\illegible{1}}\inserted{$\widehat{\unsure{\mathcal{G}}}_{0,3}^{+}$} $\simeq \widehat{\mathcal{T}}_{1,1}^{+\prime}$, \deleted{l'hom.}

\margin{30}

\scanpage{32}
\mspage{304}

a) l'extension plus fine $N_{1,1}$ de $\Gamma_{\mathbb{Q}}$ par $\widehat{\mathcal{T}}_{1,1}^{+}$, qui est donc une ext. de $N'_{1,1}$ (\unsure{censé} être connu) par le groupe $\{\pm1\}$

b) l'hom
\[
\begin{array}{|c|}
\hline
N_{1,1}\longrightarrow \mathrm{Gl}(2,\widehat{\mathbb{Z}})
\end{array}
\tag{21}
\]
de (1), et à défaut, l'hom
\unsure{\deleted{$N'_{1,1}\simeq\widehat{\mathcal{G}}_{0,3}^{+}$}}
\[
N'_{1,1}\simeq
\overset{\inserted{\unsure{\text{nouveau}}}}{\deleted{\unsure{\mathcal{G}}}\inserted{M}}_{0,3}
\longrightarrow
\mathrm{Gl}(2,\widehat{\mathbb{Z}})/\{\pm1\}.
\tag{22}
\]

Je \unsure{commencerai|commence} par la question b).

Si on dispose de l'extension $N_{1,1}$ de $\Gamma_{\mathbb{Q}}$ par $\widehat{\mathcal{T}}_{1,1}^{+}\inserted{\simeq\mathrm{Sl}(2,\mathbb{Z})^{\wedge}}$, on a un hom
\[
N_{1,1}\longrightarrow
\mathrm{Aut}_{\mathrm{lac}}
\left(\underbrace{\widehat{\mathcal{T}}_{1,1}^{+}}_{\mathrm{Sl}(2,\mathbb{Z})^{\wedge}}\right)
\tag{23}
\]
qui passe au quotient en un hom
\[
N_{1,1}\longrightarrow
\mathrm{Aut}_{\mathrm{lac}}\left(\mathrm{Sl}(2,\widehat{\mathbb{Z}})\right)
\tag{24}
\]
or on a un hom évident
\[\deleted{\unsure{(2}}athrm{Gl}(2,\widehat{\mathbb{Z}})
\longrightarrow
\mathrm{Aut}_{\mathrm{lac}}\left(\mathrm{Sl}(2,\widehat{\mathbb{Z}})\right)
\]
qui se factorise en
\[
\mathrm{Gl}(2,\widehat{\mathbb{Z}})
\longrightarrow
\mathrm{Gl}(2,\widehat{\mathbb{Z}})/\widehat{\mathbb{Z}}^{*}
\simeq \mathrm{GP}(1,\widehat{\mathbb{Z}})
\hookrightarrow
\mathrm{Aut}_{\mathrm{lac}}\left(\mathrm{Sl}(2,\widehat{\mathbb{Z}})\right)
\tag{25}
\]
et il se trouve que (24) se factorise en
\[
N_{1,1}\longrightarrow
\underset{\simeq\,\mathrm{Gl}(\unsure{2|1},\widehat{\mathbb{Z}})/\widehat{\mathbb{Z}}^{*}}
{\mathrm{GP}(1,\widehat{\mathbb{Z}})}
\hookrightarrow
\mathrm{Aut}_{\mathrm{lac}}\left(\mathrm{Sl}(2,\widehat{\mathbb{Z}})\right)
\tag{26}
\]

\margin{31}

\scanpage{33}

ce qui permet donc de reconstituer, à partir de l'extension $N_{1,1}$, un hom
\[
N_{1,1}\xrightarrow{\varphi}
\mathrm{GP}(1,\widehat{\mathbb{Z}})
\simeq
\mathrm{Gl}(2,\widehat{\mathbb{Z}})/\widehat{\mathbb{Z}}^{*},
\tag{27}
\]
un peu plus faible que l'hom (21) qu'on cherche à construire. Mais notons que l'hom
\[
\diagram{
\begin{tikzcd}[ampersand replacement=\&, column sep=small,row sep=small]
\mathrm{Gl}(2,\widehat{\mathbb{Z}})
\arrow[r,"{u\mapsto(\dot{u},\det(u))}"]
\&
\mathrm{GP}(1,\widehat{\mathbb{Z}})\arrow[d]
\&
{}\times\widehat{\mathbb{Z}}^{*}
\\
\&
\widehat{\mathbb{Z}}^{*}/\widehat{\mathbb{Z}}^{*2}
\&
\end{tikzcd}
}
\tag{28}
\]
a comme \uline{noyau} l'ens des $\lambda\in\widehat{\mathbb{Z}}^{*}$ tels que $\lambda^2=1$ -- i.e. le groupe
\[
{}_{2}\widehat{\mathbb{Z}}^{*}
\simeq (\mathbb{Z}/2\mathbb{Z})^{P}
\tag{29}
\]
\margin{$P\unsure{=|:}$ ens des nbs premiers}

Ainsi, on trouve que l'hom
\[
\deleted{\unsure{\Gamma_{\mathbb{Q}}}}\inserted{N_{1,1}}
\xrightarrow{u\mapsto(\varphi(u),\chi(u))}
\mathrm{GP}(1,\widehat{\mathbb{Z}})\times\widehat{\mathbb{Z}}^{*}
\tag{30}
\]
(déduit de (27) et de $\chi$) \unsure{\S}
\margin{c'est encore une \uline{condition} sur l'ext. $N_{1,1}$ de $\Gamma_{\mathbb{Q}}$ par $\widehat{\mathcal{T}}_{1,1}^{+}$, i.e. \unsure{de} l'op. ext. de $\Gamma_{\mathbb{Q}}$ sur $\widehat{\mathcal{T}}_{1,1}^{+}$}
\uline{se factorise par l'image de} (28), et définit donc un hom
\[
\deleted{\unsure{\Gamma_{\mathbb{Q}}}}\inserted{N_{1,1}}
\xrightarrow{\psi}
\mathrm{Gl}(2,\widehat{\mathbb{Z}})/{}_{2}\widehat{\mathbb{Z}}^{*}
\tag{31}
\]
qui est encore une version affaiblie de (21).

\margin{32}

\scanpage{34}
\mspage{305}

À vrai dire, l'homomorphisme (31) se factorise en
\[
\tag{32}
\bigl[N_{1,1}\longrightarrow\bigr]
\underset{\substack{\Vert\\N_{1,1}/\pm1}}{N'_{1,1}}
\xrightarrow{\psi'}\mathrm{Gl}(2,\widehat{\mathbb Z})/{}_{2}\widehat{\mathbb Z}^{*}
\]
et l'hom $\psi'$ peut sûrement se construire directement à partir de l'extension
$N'_{1,1}$ de $\Gamma_{\mathbb Q}$ par
$\widehat{\mathcal T}_{1,1}^{+\prime}\simeq\unsure{\widehat{\mathcal G}_{0,3}^{+}}$
\inserted{$\simeq\mathrm{Sl}(2,\widehat{\mathbb Z})'$}, tout comme tantôt $\unsure{\psi|\varphi}$ :
il suffit de vérifier que l'hom canonique
\[
\tag{33}
\underset{\substack{\unsure{\wr}\\\mathrm{Gl}(2,\widehat{\mathbb Z})/\widehat{\mathbb Z}^{*}}}{\mathrm{GP}(1,\widehat{\mathbb Z})}
\longrightarrow\mathrm{Aut}_{\mathrm{lac}}\bigl(\mathrm{Sl}(2,\widehat{\mathbb Z})'\bigr)
\]
est \inserted{encore} injectif. La question est ici
\inserted{\unsure{(en termes de th. des gr. profinis)}}
comment \deleted{relever \illegible{1}}\inserted{expliciter le relèvement de}
$\psi'$ de \unsure{(31)} en l'homom. canonique
\[
\tag{34}
\deleted{\unsure{\Gamma_{\mathbb Q}}}\quad
N'_{1,1}\longrightarrow\mathrm{Gl}(2,\widehat{\mathbb Z})'
\overset{\mathrm{dfn}}{=}\mathrm{Gl}(2,\widehat{\mathbb Z})/\pm1
\]
--- une fois cet hom obtenu, on en déduit l'ext. $N_{1,1}$ par image inverse
de la structure d'ext.
\[
\tag{35}
1\longrightarrow\{\pm1\}\unsure{\longrightarrow}\mathrm{Gl}(2,\widehat{\mathbb Z})
\longrightarrow\mathrm{Gl}(2,\widehat{\mathbb Z})'\longrightarrow1,
\]
ainsi que l'hom (21) $N_{1,1}\longrightarrow\mathrm{Gl}(2,\widehat{\mathbb Z})$.

\margin{33}

\scanpage{35}

Donc en résumé, la question essentielle semble être la détermination du relèvement (34) de (32).

Reprenons nos arguments. On se donne un sous-groupe $\Gamma$ de
$\mathrm{Autext}_{\mathrm{lac}}(\widehat{\Pi}_{0,3})$, commutant à
$\unsure{\mathfrak S_3}$, d'où une extension \deleted{\illegible{1}}
(notée $M_{0,3}$ ou $N'_{1,1}$) de $\Gamma$ par
$\unsure{\widehat{\mathcal G}_{0,3}^{+}}$, ou encore de
$\unsure{\mathfrak S_3}\times\Gamma$ par $\widehat{\Pi}_{0,3}$,
\deleted{\unsure{Utilisant} l'isom
$\unsure{\widehat{\mathcal G}_{0,3}^{+}}\simeq\mathrm{Sl}(2,\widehat{\mathbb Z})'$, on trouve un}
\inserted{d'où un} homomorphisme
\[
\tag{36}
M_{0,3}\unsure{\bigl(N'_{1,1}\bigr)|\bigl(=N'_{1,1}\bigr)}\longrightarrow
\mathrm{Aut}_{\mathrm{lac}}\left(
\underbrace{\unsure{\widehat{\mathcal G}_{0,3}^{+}}}_{\mathrm{Sl}(2,\unsure{\widehat{\mathbb Z}|\mathbb Z})'{}^{\wedge}}
\right)
\]
induisant sur $\unsure{\widehat{\mathcal G}_{0,3}^{+}}$
\deleted{\unsure{l'opération de $\Gamma$}}\inserted{l'opération}
int. de ce groupe sur lui-$\overline{\mathrm m}$.
On a par ailleurs un hom injectif
\[
\tag{37}
\mathrm{GP}(1,\widehat{\mathbb Z})\hookrightarrow
\mathrm{Aut}\bigl(\mathrm{Sl}(2,\widehat{\mathbb Z})'\bigr)
\]
et $\Gamma\subset\mathrm{Autext}_{\mathrm{lac}}(\widehat{\Pi}_{0,3})$ est
\unsure{astreint} (si ce n'est \inserted{déjà} une conséquence de l'hyp. de commutation à
$\unsure{\mathfrak S_3}$) à la condition que l'hom
\inserted{ext.}\deleted{\illegible{1} ext. \unsure{de $\Gamma$}}
\[
\tag{38}
\underset{\substack{\unsure{\wr}\\\mathrm{Sl}(2,\mathbb Z)'{}^{\wedge}}}{\unsure{\widehat{\mathcal G}_{0,3}^{+}}}
\longrightarrow\mathrm{Sl}(2,\widehat{\mathbb Z})'
\]

\margin{34}

\scanpage{36}
\mspage{306}

soit compatible aux actions ext. de $\Gamma$, sur
$\unsure{\widehat{\mathcal G}_{0,3}^{+}}$ par tautologie (du fait que $\Gamma$ est formé
d'él. de $\mathrm{Autext}_{\mathrm{lac}}(\widehat{\Pi}_{0,3})$ qui commutent à
$\unsure{\mathfrak S_3}$), sur $\mathrm{Sl}(2,\widehat{\mathbb Z})'$ via
$\widehat{\mathbb Z}^{*}/\widehat{\mathbb Z}^{*2}$ et le caractère ``cyclotomique''
\[
\tag{39}
\chi:\Gamma\longrightarrow\widehat{\mathbb Z}^{*}
\]
déduit de l'op. ext. tautologique de $\Gamma$ sur $\widehat{\Pi}_{0,3}$
(respectant la structure à lacets).

Cette ``condition de congruence'' \deleted{\unsure{signifie}}
\deleted{\unsure{implique}} \unsure{aussi|ainsi} que (36)
\inserted{définit
$\unsure{M_{0,3}|N'_{1,1}}\longrightarrow\mathrm{Aut}\,\mathrm{Sl}(2,\widehat{\mathbb Z})'$, qui}
(se factorise par (37)), en un hom \unsure{bien} déterminé
\[
\tag{40}
\varphi':M_{0,3}^{\deleted{\unsure{+}}}\longrightarrow\mathrm{GP}(1,\widehat{\mathbb Z}),
\]
d'où, en utilisant de plus ledit caractère cyclotomique, un hom
\[
\tag{41}
\deleted{\unsure{M}_{0,}}\quad M_{0,3}
\xrightarrow{\inserted{u\mapsto(\varphi'(u),\chi(u))}}
\mathrm{GP}(1,\widehat{\mathbb Z})\times\widehat{\mathbb Z}^{*}
\]
et la condition \unsure{congruentielle} nous assure en fait, plus précisément, que ce dernier
se factorise en un hom
\[
\tag{42}
M_{0,3}\xrightarrow{\psi'}\mathrm{Gl}(2,\widehat{\mathbb Z})/{}_{2}\widehat{\mathbb Z}^{*},
\]
compte tenu de l'injection

\margin{35}

\scanpage{37}
\[
\begin{aligned}
(43)\qquad
\mathrm{Gl}(2,\widehat{\mathbb Z})/{}_{2}\widehat{\mathbb Z}^{*}
&\hookrightarrow \mathrm{GP}(1,\widehat{\mathbb Z})\times\widehat{\mathbb Z}^{*}\\
u\bmod {}_{2}\widehat{\mathbb Z}^{*}
&\longmapsto (u\bmod\widehat{\mathbb Z}^{*},\det u).
\end{aligned}
\]

Considérons l'hom can. (surjectif)
\[
(44)\qquad
\underset{\substack{\mathrm{dfn}\;\Vert\\
\mathrm{Gl}(2,\widehat{\mathbb Z})/\pm1}}
{\mathrm{Gl}(2,\widehat{\mathbb Z})'}
\longrightarrow \mathrm{Gl}(2,\widehat{\mathbb Z})/{}_{2}\widehat{\mathbb Z}^{*},
\]
dont le noyau est \unsure{isom.} :
\[
(45)\qquad
\unsure{\nu|V}=(\mathbb Z/2\mathbb Z)^P/\mathrm{diag}
\qquad (P=\text{ens des nbs. premiers}).
\]

L'hom. canonique $\psi'$ (42), déduit des conditions de congruence,
définit par image inverse une extension centrale
\[
(46)\qquad
1\longrightarrow\unsure{\nu|V}\longrightarrow
\unsure{\widetilde{M}}_{0,3}\longrightarrow
M_{0,3}\longrightarrow1,
\]
et la donnée d'un hom (34)
\[
(47)\qquad
M_{0,3}\longrightarrow\mathrm{Gl}(2,\widehat{\mathbb Z})'
\]
qui relèverait (42), équivaut à \inserted{celle d'}un
scindage de cette extension (46). La question à élucider à présent :
est-il possible d'expliciter un tel \unsure{relèvement}\inserted{ canonique} de
(42), \unsure{ou} scindage de (46), en termes de la

\margin{36}

\scanpage{38}
\mspage{307}
seule donnée de $\Gamma\subset\mathrm{Autext}_{\mathrm{lac}}
(\widehat{\Pi}_{0,3})$, satisfaisant aux conditions congruentielles
énoncées plus haut ?

Notons que sur le sous-groupe
$\widehat{\unsure{\mathcal T|\mathcal C}}_{0,3}^{+}
\simeq \mathrm{Sl}(2,\unsure{\widehat{\mathbb Z}|\mathbb Z})'{}^{\wedge\deleted{\unsure{+}}}$
de $M_{0,3}$, le \deleted{\unsure{hom}}\inserted{relèvement}
cherché (47) est déjà connu, comme étant
\[
\mathrm{Sl}(2,\mathbb Z)'{}^{\wedge}\longrightarrow
\mathrm{Sl}(2,\widehat{\mathbb Z})'\hookrightarrow
\mathrm{Gl}(2,\widehat{\mathbb Z})'.
\]
En d'autres termes, le scindage cherché de (46) est déjà connu au dessus
du sous-groupe $\widehat{\unsure{\mathcal T|\mathcal C}}_{0,3}^{+}$
de $M_{0,3}$.

L'indétermination des rel. ou scindages étant a priori dans
\deleted{\illegible{1}}
\[
\mathrm{Hom}(M_{0,3},\uline{\unsure{\nu|V}}),
\]
tenant compte de la condition précédente, elle sera dans
\[
\begin{aligned}
(48)\qquad
\mathrm{Ker}\bigl(\mathrm{Hom}(M_{0,3},\uline{\unsure{\nu|V}})
&\longrightarrow
\mathrm{Hom}(\widehat{\unsure{\mathcal T|\mathcal C}}_{0,3}^{+},
\uline{\unsure{\nu|V}})\bigr)\\
&\simeq\mathrm{Hom}(\Gamma,\unsure{\nu|V})\;-
\end{aligned}
\]
si on suppose que
$\unsure{\Gamma/\Gamma_{\mathrm{ab}}}\xrightarrow{\chi}\widehat{\mathbb Z}^{*}$
est iso, l'indétermination \unsure{serait} encore dans
$\mathrm{Hom}((\widehat{\mathbb Z}^{*})_2,\unsure{\nu|V})$
--- c'est gros ! Il

\margin{37}

\scanpage{39}
faut un élément un peu \unsure{astucieux} pour justifier ce relèvement --- on
devine qu'il est lié justement à l'interprétation d'un $M_{0,3}$ comme
un $N'_{1,1}$ --- i.e. à l'introduction, dans une ``théorie profinie
\unsure{de|des} courbes rationnelles'', de la ``théorie profinie des
courbes elliptiques''.

\figure{\begin{tikzpicture}[x=1cm,y=1cm,line cap=round]
\draw[line width=.28pt] (0,0) .. controls (3.4,.16) and (7.4,.47) .. (11.65,.69)
 .. controls (11.78,.70) and (11.85,.68) .. (11.91,.68);
\end{tikzpicture}}

\uline{40} \unsure{Reformulation}
\inserted{\unsure{topologico-isotopique|topologico-isotopiques}}
(en termes des groupes fondamentaux) de la Théorie de Legendre

Pour tâter le terrain, j'ai envie d'abord de partir de la situation des
\unsure{groupes} discrets, qu'on a mieux en mains, c'est le cas de le dire !

Soit $C_{0,3}$ la catégorie des groupes (discrets) à lacets de type $0,3$
\inserted{(avec les \unsure{iso.} comme flèches)}, dont le
représentant-type est $\Pi_{0,3}$ --- cette catégorie est donc
équivalente canoniquement à celle des torseurs sous le groupe
\[
(49)\qquad
\mathrm{Aut}_{\mathrm{lac}}(\Pi_{0,3})
\simeq\unsure{\mathcal G|\mathcal S}_{0,3}
\simeq\mathrm{Gl}(2,\mathbb Z)^{\unsure{\wedge}}/\pm1
\simeq\mathrm{GP}(1,\mathbb Z)
\]
Soit $C_{1,1}^{\mathrm{ext}}$ la catégorie des groupes à lacets
\inserted{\uline{extérieurs}} de type $(1,1)$,
\deleted{\unsure{elle est}}\inserted{dont le} représentant-type est
$\Pi_{1,1}=\{x,y,l\mid[x,y]l=1\}$ ; cette catégorie est

\margin{38}

\scanpage{40}
\mspage{308}
équivalente à celle des torseurs sous le groupe
$\mathrm{Autext}_{\mathrm{lac}}(\Pi_{1,1})
\unsure{\xrightarrow{\sim}|\simeq}
\inserted{\mathcal T_{1,1}\simeq}\mathrm{Gl}(2,\mathbb Z)$
--- l'isomorphisme
\[
(50)\qquad
\mathrm{Autext}_{\mathrm{lac}}(\Pi_{1,1})\longrightarrow
\mathrm{Gl}(2,\mathbb Z)
\]
: il est obtenu en associant à tt automorphisme extérieur de
$\Pi_{1,1}$ son action sur
\[
(51)\qquad
(\Pi_{1,1})_{\mathrm{ab}}\simeq\mathbb Z^2
\]

On a donc un foncteur canonique
\[
(52)\qquad
L:\deleted{\unsure{A}}\ C_{1,1}^{\mathrm{ext}}
\longrightarrow C_{0,3}
\]
\margin{$L$, initiale de ``\uline{Legendre}''}
défini (à isom \unsure{can.} près) par la condition de rendre commutatif
le diagramme

(53)\quad
\diagram{\begin{tikzcd}[ampersand replacement=\&]
C_{1,1}^{\mathrm{ext}} \arrow[r] \arrow[d, "\wr"] \& C_{0,3} \arrow[d] \\
\mathrm{Tors}(\mathrm{Gl}(2,\mathbb Z)) \arrow[r] \&
\mathrm{Tors}(\mathrm{GP}(1,\mathbb Z))
\end{tikzcd}}

où \unsure{l'|1'}
\[
\mathrm{Tors}(\mathrm{Gl}(2,\mathbb Z))\longrightarrow
\mathrm{Tors}(\mathrm{GP}(1,\mathbb Z))
\]
est le foncteur de chgt de groupe d'\unsure{opérateurs|opérations} relatif à :
\[
(54)\qquad
\mathrm{Gl}(2,\mathbb Z)\xrightarrow{\mathrm{can}}
\mathrm{GP}(1,\mathbb Z)\simeq
\mathrm{Gl}(\unsure{1|2},\mathbb Z)/\pm1
\]

\margin{39}
